\section{Threats to Validity and Limitations}\label{sec:limitations}\label{sec:threats}

The preceding sections report several results this paper is obligated, by its
own stated methodology, never to soften: a negative topology lift on SAP
W\"urzburg, an uninformative comparison on a saturated IBM AML baseline, a
single-seed ablation, and a deployed model that is provably not the model
evaluated as the research champion. This section collects those admissions in
one place, cross-referenced rather than re-derived, and adds several threats
not yet raised elsewhere in the paper: the measured cost of naive synthetic
augmentation, a numerical-stability defect that quietly touches every headline
number in this study rather than only the serving path, and the scope
boundaries of the production system's access-control and deployment model.
Table~\ref{tab:threats-summary} summarizes all seven; the remainder of this
section explains each.

\begin{table}[t]
\centering
\caption{Threats to validity, at a glance. Every magnitude is read directly from a named artifact file, never re-derived here.}
\label{tab:threats-summary}
\resizebox{\columnwidth}{!}{%
\scriptsize
\begin{tabular}{p{2.6cm}p{4.0cm}p{2.6cm}}
\toprule
Threat & Magnitude / evidence & Disclosure \\
\midrule
Single-seed ablation, no CIs & IBM AML full-model lift $+0.00792$ is below the project's own $\sim$0.01 noise floor & \S\ref{sec:ablation-stats}; 3-seed champion partially corroborates \\
\addlinespace
TGN not bit-reproducible & BankSim $0.938\to0.845$; SAP W. $0.140\to0.098$, fixed seed, same source data, different machine & \texttt{docs/BUILD\_SPEC.md}; \S\ref{sec:system-serving} \\
\addlinespace
Batch z-score instability & 1{,}238/594{,}643 (0.21\%) BankSim rows, $|z|>10^3$--$10^8$ from catastrophic cancellation & Online engine uses Welford; batch pipeline uncorrected \\
\addlinespace
Synthetic augmentation hurts & $\Delta$PR-AUC $-0.1059$/$-0.00455$/$-0.10725$ on 3 real datasets, all 3 seeds each & Never used in the deployed bundle's training data \\
\addlinespace
SAP W. small test fold & 25 test frauds; 3-seed champion std $\pm0.0717$ on mean $0.1557$ (46\% rel.) & Flagged high-variance at every occurrence \\
\addlinespace
Deployed $\neq$ champion & Online-bundle gap $0.005$--$0.208$ PR-AUC by dataset & Table~\ref{tab:online-gap}; every PR-AUC labeled champion/online \\
\addlinespace
Coarse RBAC, single tenant & 3 roles, no maker--checker chain; no cross-org isolation tested & Named explicitly as unevaluated scope \\
\bottomrule
\end{tabular}}
\end{table}

\subsection{Statistical power: one seed is not an estimate}\label{sec:limitations-stats}
The seven-arm ablation that produces this paper's headline lift numbers
(Table~\ref{tab:ablation-lifts}) is a single training run per arm, seed 42,
with no repeated trials, no confidence intervals, and no significance test of
any kind. \S\ref{sec:ablation-stats} already applies the project's own
$\sim$0.01 PR-AUC noise-floor heuristic to this fact and concludes that IBM
AML's full-model structure lift ($+0.00792$) cannot be distinguished from zero
on the evidence available; we restate it here as what it is: not a p-value,
but an informal threshold the project's own harness code uses to decide when
a difference is worth interpreting at all. No formal bootstrap or paired test
has been run anywhere in this study to replace that heuristic. The separate
three-seed champion runs in \texttt{final\_champions.json} corroborate the
qualitative picture -- BankSim stable at $0.9364\pm0.0015$, SAP W\"urzburg
volatile at $0.1557\pm0.0717$ -- but they evaluate a different, TGN-inclusive
feature set on a (mostly) different split convention, so they are
corroboration, not a repetition of the ablation itself. SAP W\"urzburg's 25
test-split frauds are, on their own, close to a hard ceiling on what any
single-dataset PR-AUC estimate on this data can support; we treat its
negative lift as directionally credible but numerically imprecise throughout,
never as a precise point estimate. Even the split denominators behind this
result are not perfectly reconciled: \texttt{ablation.json} and
\texttt{final\_champions.json} report two slightly different partition-aware
split sizes for SAP W\"urzburg ($120{,}411$/$40{,}137$/$40{,}139$ vs.\
$120{,}341$/$40{,}114$/$40{,}116$, a $0.06\%$ difference, \S\ref{sec:protocol}),
traced to separate ETL passes but not further diagnosed.

\subsection{Two disclosed non-reproducibilities}
Two independent findings surfaced during this build's own reproduction
attempts, both load-bearing enough to affect how every TGN-inclusive number
in this paper should be read. First, re-training the TGN encoder on freshly
downloaded copies of the same source data, with the same fixed seed, does
not recover the committed champion bundle's test PR-AUC: BankSim recomputes
to $0.845$ against the bundle's recorded $0.938$, and SAP W\"urzburg to
$0.098$ against $0.140$ (\texttt{docs/BUILD\_SPEC.md}; \S\ref{sec:system-serving}).
CPU-side floating-point differences in PyTorch training compound across
epochs and propagate through a feature block carrying 1.9--42.0\% of the
champion model's SHAP importance depending on dataset. Practically, the
champion PR-AUC figures reported in this paper are this machine's own
measurement, not a constant guaranteed to reproduce on different hardware
with the same seed and code -- a property one would not expect to have to
state for a gradient-boosted tree, but must, because the TGN stage feeding
it is not deterministic across machines.

Second, and unrelated to cross-machine reproducibility, the batch feature
pipeline's expanding per-account z-score uses the textbook two-pass variance
formula ($\overline{x^2} - \bar{x}^2$), subject to catastrophic cancellation
for accounts whose historical amounts are nearly identical. Quantified on the
full BankSim dataset, this produces spurious $|z|>1{,}000$ values (observed
up to $\sim\!10^8$) on $1{,}238$ of $594{,}643$ rows ($0.21\%$), where a
stable computation returns a near-zero or undefined value. This is not
merely a serving-time footnote: the identical formula runs across all four
datasets' batch builds, so every headline ablation and champion PR-AUC in
this paper -- not only the BankSim figure it is quantified on -- was computed
with this instability present somewhere upstream, at an unmeasured rate on
the other three datasets. The online engine
(\texttt{app/services/online\_ordinary.py}) uses Welford's algorithm
instead, which does not fail this way; this was a deliberate choice to ship
correct code going forward rather than replicate a known bug for train/serve
identity, and it is the one place the offline and online paths diverge by
construction rather than by TGN's absence alone. The batch pipeline itself
remains uncorrected -- fixing it would move this paper's own headline numbers
after the fact -- so we disclose the defect precisely rather than patch it
mid-study.

\subsection{Synthetic data: barred as evidence, and measurably harmful as augmentation}
synth\_erp is this project's own generated corpus, and the project's binding
policy treats it as a stress-testing environment only, never as thesis
evidence -- precisely to avoid grading its own homework with data built to
contain the structure the paper argues for. Its largest-in-study topology
lift ($+0.4031$ PR-AUC, $246.8\%$ relative) is reported in \S\ref{sec:ablation}
under that explicit caveat and is not treated as support for this paper's
central claim.

A second, distinct question -- does adding synth\_erp rows to a \emph{real}
dataset's training set help a real model -- was tested directly
(\texttt{paper/data/augmentation.json}, three seeds, $1{,}148{,}503$ synthetic
rows appended to each real training split) and answers plainly: it does not.
Augmentation reduces test PR-AUC on every one of the three real datasets, at
every seed: BankSim $-0.1059\pm0.0011$ ($0.8669\to0.7610$), IBM AML
$-0.00455\pm0.0013$ ($0.99999\to0.99544$), and SAP W\"urzburg
$-0.1073\pm0.0417$ ($0.1366\to0.0293$). This is the more actionable of the
two synthetic-data warnings: not only is synth\_erp barred as evidence, but
naively mixing it into a real model's training data measurably degrades that
model, plausibly by shifting the decision boundary toward synthetic-specific
regularities that do not hold on real transactions -- which is why the
deployed online bundle (\S\ref{sec:system-serving}) is trained exclusively
on real data. One number in the same experiment cuts the other way and
answers a narrower question: a model trained \emph{only} on synth\_erp, with
zero real BankSim rows, still reaches PR-AUC $0.7243\pm0.0079$ scored
zero-shot against real BankSim test data (versus $0.8669\pm0.0031$
real-trained) -- evidence the generator's structural realism transfers to
some degree, a claim about generator fidelity rather than about whether
\emph{mixing} synthetic rows helps a real deployment. We report both rather
than only the one that reads more favorably.

\subsection{The deployed model is not the evaluated champion}
Every PR-AUC attributed to a ``champion'' in this paper describes an
offline, TGN-inclusive model that is not, and by construction cannot be,
served online (\S\ref{sec:system-serving}): TGN's per-account memory tensor
is never persisted, and TGN training is the non-reproducible component
above, so a production system should not be handed authority to reload it.
Table~\ref{tab:online-gap} quantifies what removing TGN costs against the
same dataset instance: a small $0.00523$ PR-AUC gap on BankSim, a larger
$0.04631$ on SAP W\"urzburg, and a substantial $0.20827$ on synth\_erp --
proportional to how much of that dataset's champion SHAP importance TGN
carried (17.5\%, 42.0\%, 37.3\%). IBM AML is excluded from that comparison
because its online bundle trained on a separate, fresh Kaggle re-download
(\texttt{ibm\_aml\_kaggle}), not the simulator instance behind its champion
figure, so no single gap number honestly describes it. The risk this creates
is real: the abstract's headline lifts ($+11.7\%$ BankSim, $+246.8\%$
synth\_erp, non-evidentiary) both describe the offline champion, and a
reader who quotes them as the deployed console's accuracy would overstate it
on every dataset except BankSim, where the gap is small.

\subsection{System scope: access control and deployment breadth}
The audit console enforces three named roles -- \texttt{erp\_clerk},
\texttt{forensic\_auditor}, \texttt{system\_administrator} -- through one
route-level guard (\texttt{app/core/deps.py:require\_role}), surfaced through
two front-end experiences: clerk-facing ERP screens that post invoices,
vendors, and payments, and an auditor-facing console that reviews alerts.
This is coarse by design: no maker--checker approval chain (a second clerk
or auditor co-signing a payment above a threshold), no per-vendor or
per-amount delegation, and no attribute-based policy beyond the three fixed
roles. The deployment is single-tenant: the scoring registry keys one
\texttt{FraudScoringService} instance per dataset, standing in for a client
or business unit in this build, but no cross-tenant data isolation,
per-tenant model versioning, or multi-organization quota behavior was
implemented or tested. Federated training across institutions (a
FedProx-style scheme letting organizations improve a shared model without
exchanging raw transaction data) was scoped only as a future direction in
the source project and was never implemented, benchmarked, or evaluated
here; nothing in this architecture should be read as evidence toward a
federated claim. Two narrower validity gaps round out the picture: the
four modeled datasets are static historical snapshots -- none model an
adversary adapting after observing what the deployed model flags, a concern
the tabular-evaluation literature increasingly treats as first-class rather
than an afterthought~\cite{jesus2022turning} -- and every arm in this study
is a gradient-boosted tree with at most one learned-graph baseline (TGN),
considerably narrower than the architecture space the broader graph-fraud
literature explores~\cite{motie2024gnn,yuan2025survey}.

\section{Conclusion}\label{sec:conclusion}

\subsection{The claim this paper actually supports}
This paper set out to answer one question honestly, in either direction:
does leakage-free graph topology add predictive value over a strong
transaction-only baseline? The answer, on the evidence gathered across four
datasets under one shared, leakage-audited protocol, is conditional rather
than universal, and we state it exactly as measured rather than rounded
toward a cleaner story. On BankSim, topology helps decisively: a
$+0.0974$ PR-AUC full-model structure lift ($11.7\%$ relative over a $0.833$
baseline), independently confirmed by a clean, amount-blind probe that still
finds $+0.1899$ PR-AUC once every amount feature is stripped from both sides
of the comparison. On the project's own synth\_erp corpus, topology helps by
the largest margin measured in this study ($+0.4031$, $246.8\%$ relative) --
but this number is explicitly barred from standing as thesis evidence by the
project's own policy, since the corpus was built to contain the structure
being tested for, and stands only as a demonstration of what topology
\emph{can} do under favorable conditions. On IBM AML, the comparison that
matters is uninformative rather than negative: an already-saturated $0.990$
baseline leaves almost no headroom, the nominal full-model lift
($+0.00792$) sits below this project's own noise floor, and the only honest
positive signal is a narrower, amount-blind probe ($+0.0905$) that isolates a
real but modest structural contribution once the dominant amount signal is
removed from both arms. On SAP W\"urzburg, topology does not help --
$-0.08288$ PR-AUC, a genuinely negative result on this project's real
general-ledger export, reported plainly rather than explained away, with the
caveat that only 25 test-split frauds make this specific number imprecise
even as its direction is corroborated by an independent, differently-built
experiment: a 46\% relative standard deviation on the 3-seed champion PR-AUC
for the same dataset (\S\ref{sec:limitations-stats}).
Four datasets, two clear directions and two qualified ones: this is the
paper's actual result, and it is deliberately not tidier than that.

We regard this qualified answer as more valuable than an inflated universal
one would have been, for a specific, falsifiable reason: this project's own
predecessor codebase claimed graph topology as ``the differentiator'' on
every dataset it touched, and that claim collapsed on inspection into $0.0$
feature importance and leakage-inflated metrics once anyone checked. A
four-for-four sweep of positive results, produced by the same team that
inherited that history, would be the least trustworthy possible outcome of
this study, not the most impressive one. The graph-based fraud detection
literature this paper positions itself within has, with some exceptions,
tended to report uniformly positive results for structural
features~\cite{motie2024gnn,yuan2025survey}; a controlled, same-architecture,
leakage-audited ablation that instead finds two positives, one null, and one
negative result -- and says so, including naming which of its own numbers
fall inside its own stated noise floor -- is offering a more falsifiable, and
therefore more useful, claim than a paper that reports success everywhere it
looked.

\subsection{A second, independent contribution: production engineering}
The paper's second contribution does not depend on the ablation's outcome at
all: a Procure-to-Pay ERP core with a forensic audit console was designed,
built, and \emph{measured}, not merely specified, against the hardening
requirements a real deployment would need. The online scoring path reuses
the research pipeline's own \texttt{WindowGraph} and \texttt{LapMemory}
classes directly rather than reimplementing them, verified by a parity test
that replays 20{,}000 real transactions through both the batch and online
computation paths with zero mismatches across all thirteen topology columns
(\S\ref{sec:system-arch}). Financial-integrity invariants -- double-entry
balance, fiscal-period close, and audit-log immutability -- are enforced as
Postgres triggers the application layer cannot bypass, not as conventions a
future migration could quietly break, and a five-layer duplicate-payment
defense was proven, not asserted, under genuine concurrency: 30 simultaneous
identical payment requests against one invoice produced exactly one
successful payment and 29 correctly rejected duplicates
(\S\ref{sec:system-integrity}). Performance and failure behavior are reported
as measured, including where the measurement is unflattering: the mixed
workload ramp clearly degrades before 500 concurrent users on a single
dev-mode process (p95 $8.18$s against a 1s target), and an initial
backup/restore drill correctly surfaced a Postgres MVCC-driven fingerprint
mismatch under concurrent writes rather than hiding it. None of this is
contingent on topology winning its ablation anywhere; it is a claim about
engineering discipline -- that a leakage-free feature computation can be
carried, unmodified, from a research pipeline into a live scoring path, and
that the resulting system's integrity guarantees, failure modes, and
performance ceiling can be demonstrated under real load rather than assumed.
We consider a genuinely deployed, tested, database-enforced system a
different and complementary kind of evidence from a PR-AUC number, and this
paper reports both rather than treating the console as an illustrative demo
of the research result.

\subsection{Future work}
Three directions follow directly from specific, named gaps this build
surfaced, rather than from general aspiration. First, and most concretely
scoped: persist the TGN's per-account GRU memory tensor as a checkpointed,
incrementally-updatable artifact, so that a live transaction can update and
reuse account memory in $O(1)$ rather than requiring a full history replay.
This is the single change that would let a true online champion -- topology
plus a learned temporal representation -- be served at all; today it is
blocked by an engineering gap (no persisted memory), not by an algorithmic
one, and closing it would also let the online/offline gap quantified in
Table~\ref{tab:online-gap} ($0.005$--$0.208$ PR-AUC by dataset) be measured
directly rather than only bounded. Second, extend multi-seed evaluation, and
ideally formal interval estimation, to the ablation itself rather than only
to the champion selection step: the seven-arm decomposition that produces
this paper's headline lift numbers remains single-seed, and work on
uncertainty-aware evaluation for graph-based fraud
models~\cite{shi2026uncertainty} points toward replacing this project's
informal $\sim$0.01 noise-floor heuristic with a proper bootstrap or paired
confidence interval, particularly for SAP W\"urzburg, where a 46\% relative
standard deviation in a related multi-seed experiment suggests the ablation's
point estimate alone should not be over-trusted. Third, resolve the batch
pipeline's z-score instability upstream rather than leaving it as a disclosed
but uncorrected defect: replacing the two-pass variance formula with
Welford's algorithm -- already done in the online engine -- would remove the
one place in this study where the offline research numbers and the online
serving numbers are known to diverge by construction, at the cost of
recomputing every headline PR-AUC this paper reports. We view that
recomputation, not the current disclosure-only stance, as the right eventual
resting state, and name it as such rather than treating disclosure as a
substitute for the fix.

The narrower, more durable claim this paper defends is not that graph
topology is or is not valuable for fraud detection in general, but that the
question can be asked honestly, with leakage structurally excluded by
construction rather than checked for after the fact, and that the answer one
gets is worth reporting exactly as measured -- including when it says no.
